\section{Paying Down the Tax}
\label{sec:paying}

If remoteness is a recurring tax, policy can lower the rate. The aim is not to
abolish distance but to stop it compounding. Six levers follow from the
evidence above.

\paragraph{1. Pool scale.} The cheapest way to be less small is to share.
Regional bulk procurement of medicines and fuel, shared shipping contracts,
regional universities and regulators, and joint diplomatic missions spread
fixed costs over larger populations. The Pacific already does some of this;
the evidence on thin institutions (\cref{sec:accel}) argues for much more,
especially in fast-moving areas such as cybersecurity, data protection and AI
regulation, where no small ministry can keep up alone.

\paragraph{2. Buy redundancy.} \nSIDSOneCable{} island states depend on one
cable or none. A second cable, a satellite back-up and pre-positioned repair
arrangements are insurance, and should be financed as insurance: by pooling
risk across many states rather than asking each tiny market to justify a
second system commercially. Tonga's new branch on the Hawaiki cable, built
with Australian and New Zealand support four years after the
\tongaOutageDays-day blackout, shows both what is possible and how long it
can take \parencite{aiffp2026tonga}.

\paragraph{3. Finance vulnerability, not income.} Many small states are
\enquote{too rich} for concessional finance by income per head while being far
more exposed than poorer countries (\cref{fig:climate}). Allocating
concessional finance and climate funds by measured vulnerability, rather than
by income alone, would match money to risk.

\paragraph{4. Bargain as a bloc.} The tuna Vessel Day Scheme shows what
coordinated small states can extract when they control something scarce
\parencite{aqorau2009}. The same logic applies to shipping and data
contracts, to landing rights for cables and satellites, and to climate
negotiations, where island coalitions have long punched above their weight.

\paragraph{5. Make talent circular.} Migration will continue, and it brings
remittances and skills (\cref{sec:people}). Policy can tilt it towards
circulation: portable pensions and qualifications, remote-work arrangements
that let diaspora professionals serve home institutions, bonded scholarships
paired with good jobs at home, and temporary labour schemes with real skill
transfer.

\paragraph{6. Build technology for small languages.} Digital public
infrastructure, open-source tools and AI models trained on small languages
are public goods with near-zero marginal cost. Funding them regionally---and
insisting that platforms serve small-language markets---lowers the newest and
fastest-compounding layer of the tax.
